\subsection{首视频逐题审计与局限}
2026 年 10 月 31 日 20:09 对 study\_05 单独完成了一次临时评分：
12 题中 9 题有效回答、3 题循环超时；2 题满分、4 题半分、6 题零分，
均分 0.3333，评委错误为 0。正式五视频统一评分另行执行，
不取两次评判中的较高者；故本节临时分与最终总表可能因评委波动而不同。
本节只审计首视频的预测、金标和调用记录，不是五视频完整错误分布，
也没有通过观看全部原视频独立核验金标。

\begin{itemize}
\item 彩纸回答在桌面，金标为最右书柜顶层抽屉；可能混淆原始存放与使用位置。
\item 树的总数回答 2，金标为 1，回答将``再做一棵''的计划用作完成依据。
\item 树干步骤回答 3，与金标一致，但引用理由仍需核实，不能仅凭数值匹配
证明推理过程正确。
\item 书包回答桌面，金标为桌子左侧，却获临时半分；说明自动评委可能偏宽。
\item 装饰物答纸苹果，金标是黏土苹果，对象类别部分正确但材质错误。
\item 叶子步骤、是否只关心学习和剪刀次数三题循环超时，没有有效答案；
它们不是评委故障。
\end{itemize}

首视频 QA 记录了 60 次逻辑模型调用、累计 3,994,009 输入 tokens，
单次输入峰值 147,951 tokens。长上下文与循环超时并存，
支持优先检查检索证据累积和迭代效率；但没有严格消融，不能认定
所有超时均由上下文大小导致。下一步可冻结同一份记忆，
仅切换已有有界上下文机制，分别比较答案质量、覆盖、输入量与耗时。

\subsection{另两视频的补充审计}
23:10 的临时评分中，gym\_03 为 8 题满分、3 题零分，均分 0.7273；
其中一题超时，另两题分别把果汁的中层位置答成底层、
把护腰的顶层位置答成中层。训练方式、热身道具和教练职业等回答与金标匹配。
bedroom\_04 为 7 题满分、4 题零分，均分 0.6364；
其中两题超时，另两题是过去偏好的演员答错，以及将 Mia 的薯片偏好答成巧克力。
这些例子提示应继续检查空间细节、人物--偏好绑定和时序转换，
但不构成已定位到某个具体模块的因果证据。
两视频 QA 输入峰值分别为 131,454 和 144,546 tokens，
说明长检索上下文现象不只存在于首视频。
